\chapter{Autocallables}\label{ch:dv:autocallables}

En janvier 2024, les indices CSI 500 et CSI 1000 d'actions chinoises continentales ont fortement chuté. Des notes structurées connues sous le nom de
\omterm{def:dv:autocallables:snowball}{snowballs}, avec des niveaux de knock-in compris entre 70~\% et 80~\% de leur niveau de départ, ont été activées à mesure que les indices baissaient. Les
notes versaient des coupons de 10~\% à 20~\% par an tant que les indices restaient dans leur corridor. Une fois activées, elles se sont transformées en positions qui perdaient avec l'indice. Les dealers qui les avaient émises détenaient
des couvertures qui changeaient brutalement à ces niveaux. Le même type de produit, l'\omterm{def:dv:autocallables:autocallable}{autocallable}, est vendu à des investisseurs particuliers en
Corée et au Japon. Les investisseurs coréens détenaient 19.3~billions de wons de notes liées au Hang Seng China Enterprises Index en
novembre 2023, dont 10.2~billions à échéance au premier semestre 2024~; au 22 janvier 2024, celles vendues par cinq banques commerciales avaient
déjà perdu 229.6~milliards de wons, et ces pertes ont fait l'objet d'études à part entière. Ce chapitre décompose le produit. Il traite du payoff et du modèle qu'il exige, de ses risques pour
l'émetteur, les flux de couverture qu'il engendre et ce que ces flux ont provoqué sur les marchés qu'ils couvrent.

\section{Anatomie du payoff}

\begin{definition}[Autocallable]\label{def:dv:autocallables:autocallable}
Un \emph{autocallable}\index{autocallable} est une note qui se rembourse par anticipation, au pair majoré d'un coupon, à la première
date d'observation à laquelle le sous-jacent se situe au niveau ou au-dessus d'un niveau de déclenchement. Si la note n'est jamais rappelée, elle rembourse le pair à l'échéance, sauf si une
\omterm{def:dv:autocallables:protection}{barrière de protection} a été franchie, auquel cas elle rembourse le pair multiplié par la performance du sous-jacent.
\end{definition}

\begin{definition}[Seuil d'autocall]\label{def:dv:autocallables:trigger}
Le \emph{seuil d'autocall}\index{seuil d'autocall} est le niveau, généralement de 100~\% du fixing initial et parfois
décroissant dans le temps, à partir duquel la note est remboursée par anticipation à une date d'observation.
\end{definition}

\begin{definition}[Barrière de coupon]\label{def:dv:autocallables:couponbarrier}
La \emph{barrière de coupon}\index{barrière de coupon} d'un \omterm{def:dv:autocallables:autocallable}{autocallable} de type Phoenix est le niveau à partir duquel un coupon périodique est
versé à une date d'observation, que la note soit rappelée ou non.
\end{definition}

\begin{definition}[Coupon mémoire]\label{def:dv:autocallables:memory}
Un \emph{coupon mémoire}\index{coupon mémoire} est un coupon de type Phoenix qui, s'il n'a pas été versé parce que le sous-jacent se trouvait sous la
\omterm{def:dv:autocallables:couponbarrier}{barrière de coupon}, est versé plus tard à la première date d'observation à laquelle la barrière est de nouveau atteinte.
\end{definition}

\begin{definition}[Barrière de protection]\label{def:dv:autocallables:protection}
La \emph{barrière de protection}\index{barrière de protection} (niveau de knock-in) est le niveau en dessous duquel le capital de l'investisseur
n'est plus protégé. Elle n'est observée qu'à l'échéance (européenne), ou en continu ou quotidiennement (américaine), et elle est fixée bien en dessous
du fixing initial~: 60~\% dans l'exemple du chapitre, 70~\% à 80~\% dans les \omterm{def:dv:autocallables:snowball}{snowballs} chinois de 2024.
\end{definition}

Décomposé, un \omterm{def:dv:autocallables:autocallable}{autocallable} est une obligation plus un ensemble d'options, et l'investisseur est vendeur de la plupart d'entre elles. L'investisseur a vendu
un put down-and-in dont le strike est fixé au niveau initial (le knock-in du chapitre~15)~: sous la \omterm{def:dv:autocallables:protection}{barrière de protection}, la note perd avec le
sous-jacent. L'investisseur a acheté une série d'\omterm{def:dv:barriers-and-digitals:digital}{options digitales} qui versent les coupons, et vendu la hausse au-delà du pair, car
la note est rappelée lorsque le sous-jacent monte. Le coupon est financé par la prime du put de knock-in et par le
rappel anticipé, qui raccourcit la durée de vie de la note précisément dans les scénarios où l'investisseur aurait voulu la conserver.

\begin{definition}[Snowball]\label{def:dv:autocallables:snowball}
Un \emph{snowball}\index{snowball} est un \omterm{def:dv:autocallables:autocallable}{autocallable}, répandu en Chine, dont le coupon s'accumule à un taux annuel et est versé en une
seule fois lorsque la note est rappelée par anticipation à une observation mensuelle, ou à l'échéance si elle n'a été ni rappelée ni activée~;
le knock-in est surveillé quotidiennement, et après un knock-in l'investisseur supporte à l'échéance la perte du sous-jacent.
\end{definition}

L'exemple de travail du chapitre est un Phoenix de trois ans sur l'indice de la \omterm{def:dv:implied-volatility-and-its-surface:surface}{surface de volatilité} du chapitre~9. Il est observé trimestriellement et
rappelable à partir du premier trimestre à 100~\%. Son \omterm{def:dv:autocallables:memory}{coupon mémoire} est versé à partir de 70~\%, et sa \omterm{def:dv:autocallables:protection}{barrière de protection} est observée
à l'échéance, à 60~\%. Les taux sont nuls, de sorte que la totalité du coupon est financée par des options.

\begin{example}[Où va le Phoenix]\label{ex:dv:autocallables:outcomes}
Sous la \omterm{def:dv:local-volatility:model}{volatilité locale} du chapitre~9, la note est rappelée à la première observation avec une probabilité de 57.6~\%, et dans la première
année avec une probabilité de 80.5~\%. Elle atteint l'échéance sans activation avec une probabilité de 5.7~\%, et est activée avec une probabilité de 6.3~\%
(voir la \cref{fig:dv:autocallables:outcomes}). Le coupon trimestriel équitable, celui qui valorise la note au pair, est de 1.55~\%, soit 6.18~\% par
an. Sans coupons, la note vaut 95.99, et chaque point de coupon trimestriel ajoute 2.60.
\end{example}

\begin{omfigure}
\begin{tikzpicture}
\begin{axis}[width=12cm,height=5cm,ybar,bar width=10pt,
  xlabel={issue},ylabel={probabilité (\%)},
  tick label style={font=\small},label style={font=\small},
  ymin=0,ymax=68,xmin=-0.7,xmax=12.7,
  xtick=data,xticklabels={3M,6M,9M,12M,15M,18M,21M,24M,27M,30M,33M,éch.,KI},
  x tick label style={font=\scriptsize},
  nodes near coords,nodes near coords style={font=\tiny,/pgf/number format/fixed,/pgf/number format/precision=1},
  grid=major,grid style={gray!25},table/col sep=comma]
\addplot[omThm,fill=omThm!40] table[x=i,y=prob]{figdata/derivatives/18-autocallables/outcomes.csv};
\end{axis}
\end{tikzpicture}

\omcaption{Le sort d'un Phoenix trimestriel de trois ans (\omterm{def:dv:autocallables:trigger}{seuil d'autocall} 100~\%, \omterm{def:dv:autocallables:couponbarrier}{barrière de coupon} 70~\%, protection 60~\% à l'échéance) sous
la \omterm{def:dv:local-volatility:model}{volatilité locale} du chapitre~9~: rappelé à chaque observation, remboursé au pair à l'échéance, ou activé. La plupart des notes sont rappelées dans
la première année~; la perte de l'investisseur se trouve dans la dernière barre. Données~: le tutoriel.}\label{fig:dv:autocallables:outcomes}
\end{omfigure}

\section{Valorisation~: quel modèle}

Le payoff dépend de la trajectoire à une douzaine de dates et du smile à chacune d'elles. Les coupons sont des \omterm{def:dv:barriers-and-digitals:digital}{options digitales}, donc ils sont
valorisés par le skew (chapitre~15). Le put de knock-in est une barrière très hors de la monnaie, valorisée par la partie basse du smile à
longues échéances. Et le rappel anticipé est une suite de digitales à terme. Une volatilité plate se trompe sur les trois dans le même
sens.

\begin{example}[Volatilité plate face au smile]\label{ex:dv:autocallables:flat}
Valorisé à une volatilité plate de 20.0~\%, soit la volatilité à trois ans à la monnaie, le Phoenix au coupon équitable en \omterm{def:dv:local-volatility:model}{volatilité locale} vaut 102.78
au lieu de 100. Le coupon équitable à volatilité plate est de 3.14~\% par an, soit la moitié des 6.18~\% en \omterm{def:dv:local-volatility:model}{volatilité locale}. Le skew rend
le put de knock-in, que l'investisseur a vendu, plus valorisé, et la partie courte, plus pentue, du smile réduit la probabilité d'un
rappel anticipé. Un desk qui cotait à volatilité plate paierait la moitié du coupon qu'il pourrait se permettre. Un concurrent utilisant le smile
prendrait chaque transaction.
\end{example}

La \omterm{def:dv:local-volatility:model}{volatilité locale} est le point de départ naturel~: elle recale toutes les vanilles, et la valeur de la note dépend surtout des
distributions terminales aux dates d'observation. Ce n'est pas le point d'arrivée. Les coupons et les rappels sont des digitales à terme, dont le prix
dépend du skew forward que la \omterm{def:dv:local-volatility:model}{volatilité locale} aplatit (chapitre~16). Le knock-in, surveillé quotidiennement dans de nombreux produits, dépend de la dynamique du
smile près de la barrière. Un desk prudent valorise avec un modèle à volatilité locale-stochastique (chapitre~20), le confronte à la \omterm{def:dv:local-volatility:model}{volatilité locale}, et détient une
réserve pour la différence. Il ajoute une réserve pour risque de saut ou de gap pour les barrières surveillées quotidiennement (chapitre~15). Les notes multi-sous-jacents ajoutent la corrélation, et le \omterm{def:dv:multi-asset-options:skew}{skew de corrélation} du chapitre~17.

\section{Profil de risque~: vega, skew, dividendes, corrélation}

\begin{example}[Sensibilités du Phoenix]\label{ex:dv:autocallables:greeks}
Au coupon équitable en \omterm{def:dv:local-volatility:model}{volatilité locale}, une hausse d'un point de volatilité sur toute la surface fait baisser la valeur de la note de 0.35 pour 100.
Par tranche d'échéance, la variation est de $-0.16$ à un an, $-0.05$ à deux ans et $-0.14$ à trois
ans (voir la \cref{fig:dv:autocallables:vega}). Un skew plus pentu (corrélation SSVI de $-0.6$ à $-0.7$) coûte 0.24. Un rendement de dividende de 1~\% par
an, la surface étant inchangée en moneyness forward, coûte 0.20. Entre quatre points à la baisse et six points à la hausse, la valeur
baisse presque linéairement, de 101.33 à 97.90.
\end{example}

L'investisseur est vendeur de volatilité, vendeur de skew et vendeur de dividendes, et l'émetteur détient la position inverse. L'émetteur est acheteur de \omterm{def:dv:greeks-and-the-hedging-pnl:greeks}{vega},
concentré sur les échéances courtes, où se trouvent les premières observations, et à maturité, où vit le put de knock-in. Il est acheteur de
skew et de dividendes~: une baisse des dividendes attendus fait monter le forward et réduit la valeur du put de knock-in qu'il détient. Aucune
de ces expositions ne peut rester ouverte. Le desk revend la volatilité et le skew au marché, au moyen d'options cotées et de \omterm{def:dv:variance-swaps-and-volatility-derivatives:vs}{swaps de variance}, et il
revend l'exposition aux dividendes, au moyen de contrats à terme sur dividendes et de swaps (chapitre~5). Un marché où de nombreux émetteurs vendent les mêmes notes sur les
mêmes indices dispose donc d'un offreur structurel de volatilité à long terme et d'un vendeur structurel de dividendes. Cette offre devrait
faire baisser le prix de la volatilité à long terme et des \omterm{def:dv:dividends-borrow-and-forwards:implied}{dividendes implicites}. Le chapitre ne mesure pas l'ampleur de cet effet.

\begin{omfigure}
\begin{tikzpicture}
\begin{groupplot}[group style={group size=2 by 1,horizontal sep=1.6cm},
  width=6.6cm,height=5cm,
  tick label style={font=\small},label style={font=\small},grid=major,grid style={gray!25},table/col sep=comma]
\nextgroupplot[ybar,bar width=12pt,xlabel={variation de volatilité},ylabel={variation de valeur (pour 100)},ymin=-0.4,ymax=0,xmin=-0.6,xmax=3.6,
  xtick=data,xticklabels={1A,2A,3A,total},
  title={\small vega par échéance, +1 point}]
\addplot[omThm,fill=omThm!40] table[x=i,y=vega]{figdata/derivatives/18-autocallables/vega.csv};
\nextgroupplot[xlabel={décalage parallèle de la volatilité (points)},ylabel={valeur de la note (pour 100)},xmin=-5,xmax=7,ymin=97.5,ymax=101.6,
  title={\small valeur en fonction de la volatilité}]
\addplot[omDef,very thick,mark=*,mark size=1.6pt] table[x=shift,y=price]{figdata/derivatives/18-autocallables/price_vol.csv};
\end{groupplot}
\end{tikzpicture}

\omcaption{Le Phoenix à son coupon équitable. À gauche~: la variation de valeur pour un point de volatilité ajouté autour de chaque pilier de la
surface, et sur l'ensemble. À droite~: valeur en fonction d'un décalage parallèle de la surface. L'investisseur est vendeur de \omterm{def:dv:greeks-and-the-hedging-pnl:greeks}{vega}, et l'émetteur acheteur.
Données~: le tutoriel.}\label{fig:dv:autocallables:vega}
\end{omfigure}

Les \omterm{def:dv:autocallables:autocallable}{autocallables} sont écrits sur des indices uniques et sur des paniers. Les worst-of sur trois ou quatre actions versent des coupons plus élevés parce que
l'investisseur vend davantage~: un put de knock-in sur le worst-of, dont la valeur augmente lorsque la corrélation baisse (chapitre~17).

\begin{example}[La même note sur le worst of de trois]\label{ex:dv:autocallables:worstof}
Appliquer le Phoenix au worst of de trois sous-jacents, chacun avec la \omterm{def:dv:local-volatility:model}{volatilité locale} de l'indice. Le coupon équitable vaut 11.5~\% par an, pour une
corrélation de 0.5, et 10.4~\% pour une corrélation de 0.7. La probabilité de knock-in vaut 18.1~\% et 14.6~\%. Vingt points de corrélation valent un
point de coupon annuel, et l'émetteur qui les a vendus est vendeur de corrélation~: il gagne si la corrélation monte.
\end{example}

\begin{omfigure}
\begin{tikzpicture}
\begin{axis}[width=10.5cm,height=5cm,ybar,bar width=22pt,
  xlabel={modèle et sous-jacent},ylabel={coupon équitable (\% par an)},
  tick label style={font=\small},label style={font=\small},
  ymin=0,ymax=13,xmin=-0.6,xmax=3.6,
  xtick=data,xticklabels={{vol plate},{vol locale},{worst-of 0.7},{worst-of 0.5}},
  nodes near coords,nodes near coords style={font=\footnotesize,/pgf/number format/fixed,/pgf/number format/precision=1},
  grid=major,grid style={gray!25},table/col sep=comma]
\addplot[omDef,fill=omDef!35] table[x=i,y=annual]{figdata/derivatives/18-autocallables/coupons.csv};
\end{axis}
\end{tikzpicture}

\omcaption{Le coupon annuel équitable du même Phoenix à trois ans~: à volatilité plate, sous \omterm{def:dv:local-volatility:model}{volatilité locale}, et sur le
worst of de trois sous-jacents pour des corrélations de 0.7 et 0.5. Le smile double le coupon, et le worst-of le double presque à nouveau. Données~:
le tutoriel.}\label{fig:dv:autocallables:coupons}
\end{omfigure}

\section{Flux de couverture de l'émetteur}

L'émetteur couvre le delta de la note dans le sous-jacent, et le delta d'un \omterm{def:dv:autocallables:autocallable}{autocallable} présente deux caractéristiques qui transforment
la couverture en flux de marché.

La première concerne le sens du flux. La valeur de la note pour l'investisseur monte avec le sous-jacent, si bien que l'émetteur, vendeur de la note,
détient une position acheteuse dans le sous-jacent. À mesure que le sous-jacent baisse vers la \omterm{def:dv:autocallables:protection}{barrière de protection}, le delta de la note augmente, car le
put de knock-in détenu par l'émetteur gagne rapidement de la valeur à cet endroit, et l'émetteur achète. À mesure qu'il monte vers le seuil de déclenchement, le delta de la note
diminue, et l'émetteur vend. Hors de la barrière, la couverture de l'émetteur achète les replis, vend les rebonds et amortit
la volatilité. C'est l'effet de gamma des dealers du livre~1, chapitre~26, et la raison pour laquelle les émetteurs sont acheteurs de volatilité~: ils détiennent
le gamma.

La seconde concerne la falaise au niveau d'une barrière surveillée quotidiennement. Juste au-dessus du niveau de knock-in, la note vaut le pair augmenté de son coupon si le
niveau tient et beaucoup moins s'il casse, si bien que son delta est très grand. Au knock-in, le paiement devient une simple position acheteuse dans le
sous-jacent, avec un delta d'environ un. L'émetteur doit vendre la différence, immédiatement, sur un marché en baisse. Lim et Choi ont décrit
ce schéma pour les titres coréens indexés sur actions~: le delta augmente sans cesse à mesure que le cours se rapproche du niveau de knock-in, et doit
être réduit dès que ce niveau est touché.

\begin{example}[La falaise du knock-in]\label{ex:dv:autocallables:cliff}
Un \omterm{def:dv:autocallables:snowball}{snowball} de deux ans est observé mensuellement pour le knock-out à 103~\% à partir du troisième mois, avec un knock-in quotidien à 75~\% et un coupon de
15~\% par an. En \omterm{def:dv:local-volatility:model}{volatilité locale}, il vaut 100.44 (erreur type 0.06), et sa probabilité de knock-in est de 15.6~\%. Un mois
avant l'échéance, à 25~\% de volatilité, avec le coupon plein de 30~\% dû si la barrière tient et le spot 1~\% au-dessus de la barrière, le
delta de l'émetteur est de 7.63 par unité de notionnel. Une baisse de 1~\% fait entrer la note en knock-in, et le delta tombe à 1.00. Pour 100 millions de notes, l'émetteur
vend 497 millions du sous-jacent, mesurés au niveau de la barrière. Avec trois mois restants, le chiffre est de 254 millions, et avec un
an restant 97 millions (la \cref{fig:dv:autocallables:cliff}).
\end{example}

\begin{omfigure}
\begin{tikzpicture}
\begin{axis}[width=10.5cm,height=5cm,
  xlabel={sous-jacent (\% du fixing initial)},ylabel={delta de l'émetteur (par unité de notionnel)},
  tick label style={font=\small},label style={font=\small},
  xmin=60,xmax=95,ymin=0,ymax=8.5,grid=major,grid style={gray!25},
  legend columns=2,legend style={font=\footnotesize,at={(0.5,-0.3)},anchor=north,draw=none,fill=none,/tikz/every even column/.append style={column sep=8pt}},
  table/col sep=comma]
\addplot[omThm,very thick] table[x=s,y=delta]{figdata/derivatives/18-autocallables/cliff_before.csv};
\addplot[omDef,very thick,dashed] table[x=s,y=delta]{figdata/derivatives/18-autocallables/cliff_after.csv};
\draw[gray,thick] (axis cs:75,0) -- (axis cs:75,8.5);
\node[font=\footnotesize,anchor=west] at (axis cs:76,7.9) {knock-in à 75~\%};
\legend{avant le knock-in,après le knock-in}
\end{axis}
\end{tikzpicture}

\omcaption{Le delta de l'émetteur sur un \omterm{def:dv:autocallables:snowball}{snowball} un mois avant l'échéance (knock-in à 75~\% surveillé quotidiennement, coupon de 30~\% dû si la
barrière tient, volatilité de 25~\%). L'émetteur achète à mesure que le sous-jacent baisse vers la barrière et, au knock-in, doit vendre
l'essentiel de ce qu'il a acheté. Données~: le tutoriel.}\label{fig:dv:autocallables:cliff}
\end{omfigure}

La reconstitution par Fang de l'épisode chinois des \omterm{def:dv:autocallables:snowball}{snowballs} évalue la couverture libérée dans la zone de barrière à environ 0.6 fois le notionnel
avec douze mois à l'échéance, et à environ 6 fois avec un mois, soit un ordre de grandeur comparable à celui de l'exemple. La taille du flux dépend de
trois facteurs qu'aucun émetteur seul ne contrôle~: le montant de notionnel en circulation, la concentration des niveaux de knock-in et celle des
échéances. Une douzaine d'émetteurs vendant des notes similaires aux mêmes mois, avec des barrières aux mêmes pourcentages ronds, crée une seule barrière
de très grande taille.

\section{Impact documenté sur le marché}

La structure du risque ne fait pas de doute. Les éléments sur son effet sur les marchés se multiplient, et il faut les lire en tenant compte de leurs
limites.

\begin{itemize}
\item \textbf{Les knock-in des \omterm{def:dv:autocallables:snowball}{snowballs} de 2024.} L'étude de cas de Yin sur la «~tempête de knock-in~» de janvier 2024 cite, parmi les causes, des échéances de produits concentrées
et des niveaux de knock-in regroupés, et, parmi les effets, la pression de couverture des institutionnels ainsi que la volatilité du marché.
Fang constate que l'élargissement de la base et l'activité de transaction sur les contrats à terme CSI 500 et CSI 1000, les indices les plus souvent référencés par les \omterm{def:dv:autocallables:snowball}{snowballs},
ont dépassé, autour de l'épisode, ceux des contrats à terme sur indices témoins. Il les décrit comme des résultats de cohérence, non comme une identification
causale.
\item \textbf{Titres indexés sur actions coréens.} Les études coréennes sur le delta de knock-in de l'émetteur ont abouti à une recommandation de politique~: dimensionner
les émissions sur un même sous-jacent au regard de son volume de transactions et du delta maximal que peut atteindre la couverture. Les pertes sur les notes liées au
Hang Seng China Enterprises Index font l'objet d'études empiriques de tarification.
\item \textbf{Où l'effet est faible.} Fang note que le même cadre de suivi, à savoir le notionnel en circulation multiplié par la concentration des barrières,
par la concentration des échéances et par le gamma des dealers, se transpose à d'autres marchés, mais pas son ampleur. Les notes de détail japonaises
sont plus souvent adossées à un seul sous-jacent, avec des barrières plus éloignées hors de la monnaie.
\end{itemize}

Pour un desk, les conséquences pratiques sont au nombre de trois. Connaître la concentration des barrières du marché aussi bien que la sienne~: le flux à un niveau
est celui de tous. Valoriser la falaise~: un knock-in à observation quotidienne près de l'échéance porte un \omterm{def:dv:barriers-and-digitals:gap}{risque de gap} (chapitre~15) et un coût de liquidité que le modèle
ne voit pas. Surveiller aussi le côté de l'investisseur~: le produit qui verse le coupon le plus élevé est celui dans lequel l'investisseur a vendu le
plus, et les pertes lorsque la barrière cède retombent sur l'investisseur.

\section{Tutoriel~: valorisation et couverture d'un autocallable}\label{tut:dv:autocallables:tut}

\begin{tutorial}
\textbf{Objectif.} Valoriser un Phoenix à trois ans sous \omterm{def:dv:local-volatility:model}{volatilité locale}, résoudre son coupon équitable, calculer son \omterm{def:dv:greeks-and-the-hedging-pnl:greeks}{vega} par tranche d'échéance et ses
sensibilités au skew, au dividende et à la corrélation, et mesurer la couverture de l'émetteur à travers un knock-in quotidien. \textbf{État final~:} les quatre
figures et les chiffres du problème du week-end.
\begin{tutsteps}
\item \textbf{La term sheet} sous forme de données~:
\omcode{code/firm/autocall/firm_autocall.py}{20}{34}{Term sheet de l'autocallable.}\label{lst:dv:autocallables:termsheet}
\item \textbf{Flux de trésorerie trajectoire par trajectoire}~: coupons à mémoire, rappels anticipés et remboursement à l'échéance avec le knock-in~:
\omcode{code/firm/autocall/firm_autocall.py}{63}{95}{Flux de trésorerie de l'autocallable sur les trajectoires simulées.}\label{lst:dv:autocallables:cashflows}
\item \textbf{La couverture à travers le knock-in}, à partir de la valeur de la note juste au-dessus de la barrière et juste après elle~:
\omcode{code/firm/autocall/firm_autocall.py}{128}{141}{Le sous-jacent vendu au knock-in.}\label{lst:dv:autocallables:cliff}
\item \textbf{Exécuter} \texttt{dv\_autocall.sensitivities()}, \texttt{worst\_of()}, \texttt{snowball\_price()}, \texttt{cliff()} et
\texttt{fig\_autocall.py}. Le coupon équitable repose sur la linéarité du prix en fonction du coupon~: deux simulations à nombres
aléatoires communs permettent de le résoudre.
\end{tutsteps}
\textbf{Pour aller plus loin.} Rendre la \omterm{def:dv:autocallables:protection}{barrière de protection} à observation quotidienne et recalculer le coupon équitable~; ajouter un seuil de rappel dégressif (100~\%,
95~\%, 90~\% par année)~; valoriser le Phoenix sous le \omterm{def:dv:stochastic-volatility:heston}{modèle de Heston} du chapitre~10 et comparer.
\end{tutorial}

\section{Construction~: le pricer d'autocallable}\label{bld:dv:autocallables:autocall}

\begin{build}
\bfield{Rôle} Le pricer de produits structurés de la firme miniature~: les term sheets sous forme de données, la valorisation par Monte-Carlo pour toute fonction de volatilité et tout nombre de sous-jacents,
la résolution du coupon équitable et les outils d'analyse de couverture qu'utilisent le desk de produits structurés du chapitre~19 et le livre exotique
du chapitre~27. \bfield{Interface}
\texttt{TermSheet(obs\_times, trigger, coupon, coupon\_barrier, memory, protection, ki\_daily, snowball\_rate,
first\_call)}~; \texttt{simulate(sigma\_fn, s0, t\_end, r, q, n\_paths, seed, steps\_per\_year, corr)}~; \texttt{cashflows(ts, times,
perf, r)}~; \texttt{price(ts, sigma\_fn, r, q, n\_paths, seed, corr, n\_assets) -> price, se, call\_probs, ki\_prob}~; \texttt{snowball\_tail}~;
\texttt{hedge\_across\_ki(ki, tau, vol, final\_coupon, notional)}.
\bfield{Règles} \omterm{def:dv:greeks-and-the-hedging-pnl:greeks}{Grecques} calculées par bumping avec nombres aléatoires communs~; coupon équitable déduit de la linéarité du prix en fonction du coupon~; un knock-in quotidien
valorisé sur trajectoires quotidiennes ou par \omterm{def:dv:barriers-and-digitals:shift}{décalage de barrière}, jamais en continu~; chaque prix accompagné de son erreur type.
\bfield{Tests d'acceptation} \texttt{code/firm/autocall/tests/}~: sans volatilité, la note est rappelée immédiatement~; avec un seuil jamais atteint
et sans knock-in, elle verse le pair plus tous les coupons~; un suivi quotidien fait passer en knock-in plus souvent qu'un suivi à l'échéance~; un worst-of d'actifs parfaitement corrélés se réduit à la note mono-actif~;
la forme fermée du \omterm{def:dv:autocallables:snowball}{snowball} près de l'échéance coïncide avec une simulation~; la falaise est maximale près de
l'échéance.
\bfield{Pour aller plus loin} Seuils de rappel dégressifs et \omterm{def:dv:autocallables:snowball}{snowballs} worst-of~; volatilité localement stochastique (chapitre~20)~; \omterm{def:dv:greeks-and-the-hedging-pnl:greeks}{grecques} pathwise et \omterm{def:dv:greeks-and-the-hedging-pnl:greeks}{grecques} par rapport de vraisemblance
pour les composantes digitales (chapitre~23).
\end{build}

\begin{omsources}
\item P.~Guillaume, ``Autocallable structured products'', \emph{Journal of Derivatives} 22(3) (2015) 73--94.
\item H.~Lim et Y.~Choi, ``Knock-in and stocks market effect due to ELS issuance and hedging'', \emph{Journal of Derivatives and
Quantitative Studies} 23 (2015) 289--321.
\item W.~Fang, ``China's autocallable cycle: dealer hedging flows, market impact, and lessons for Japan's structured-note market'',
SSRN 7346463 (2026).
\item Yin, ``Case study on the knock-in storm of snowball products'', \emph{Journal of Global Economy, Business and Finance} 8(3) (2026)
26--33.
\item Kim, Park et Moon, ``Markov regime-switching in pricing equity-linked securities: an empirical study for losses in HSCEI-linked
products'', \emph{Finance Research Letters} 76 (2025) 106929.
\end{omsources}

\section{Exercices}

\begin{exercise}[$\star$]\label{exo:dv:autocallables:1}
Lister les options incorporées au Phoenix du chapitre et indiquer, pour chacune, si l'investisseur est acheteur ou vendeur.
\end{exercise}

\begin{exercise}[$\star$]\label{exo:dv:autocallables:2}
Expliquer pourquoi le coupon équitable sous volatilité plate est inférieur à celui obtenu sous \omterm{def:dv:local-volatility:model}{volatilité locale}.
\end{exercise}

\begin{exercise}[$\star$]\label{exo:dv:autocallables:3}
Le prix est linéaire en fonction du coupon. À partir des chiffres du chapitre (95.99 sans coupons, 2.60 par point), calculer le coupon trimestriel
qui valorise la note à 98, en laissant à l'émetteur deux points de marge.
\end{exercise}

\begin{exercise}[$\star\star$]\label{exo:dv:autocallables:4}
Expliquer pourquoi un émetteur d'\omterm{def:dv:autocallables:autocallable}{autocallables} est long en dividendes, et ce qu'il fait pour y remédier.
\end{exercise}

\begin{exercise}[$\star\star$]\label{exo:dv:autocallables:5}
Expliquer le sens du flux de couverture de l'émetteur à mesure que le sous-jacent baisse de 100~\% vers un knock-in quotidien, et au moment du knock-in.
\end{exercise}

\begin{exercise}[$\star\star$]\label{exo:dv:autocallables:6}
Expliquer pourquoi la falaise croît à mesure que l'échéance approche.
\end{exercise}

\begin{exercise}[$\star\star\star$]\label{exo:dv:autocallables:7}
\emph{Programmation.} Rendre la \omterm{def:dv:autocallables:protection}{barrière de protection} du Phoenix à observation quotidienne et recalculer le coupon équitable sous \omterm{def:dv:local-volatility:model}{volatilité locale}. Expliquer le
changement.
\end{exercise}

\begin{exercise}[$\star\star\star$]\label{exo:dv:autocallables:8}
\emph{Trouver l'erreur.} «~Notre livre d'\omterm{def:dv:autocallables:autocallable}{autocallables} est couvert en delta et en \omterm{def:dv:greeks-and-the-hedging-pnl:greeks}{vega}, si bien qu'une chute sous le niveau de knock-in de nos notes
n'est qu'un non-événement pour le desk.~»
\end{exercise}

\section{Problème~: La falaise du knock-in}

\begin{problem}[{problème du week-end --- combien l'émetteur vend à la barrière}]\label{pb:dv:autocallables:1}
Un émetteur a vendu des \omterm{def:dv:autocallables:snowball}{snowballs} sur un indice~: knock-out mensuel à 103~\% à partir du troisième mois, knock-in quotidien à 75~\%, 15~\% par an, sur deux
ans, 100 millions de notionnel.

\medskip\noindent\textbf{Partie I --- Le produit.}
\begin{enumerate}
\item Pricer le \omterm{def:dv:autocallables:snowball}{snowball} sous la \omterm{def:dv:local-volatility:model}{volatilité locale} du chapitre~9 et donner sa probabilité de knock-in.
\item Quelles options l'investisseur a-t-il vendues~?
\item Que reçoit l'investisseur dans chacun des trois cas (knock-out, ni l'un ni l'autre, knock-in)~?
\item Pourquoi le knock-in surveillé quotidiennement est-il un risque plus difficile à gérer qu'un risque observé à l'échéance~?
\item Comparer avec le Phoenix vu dans le chapitre~: lequel rapporte le plus, et pourquoi~?
\end{enumerate}

\medskip\noindent\textbf{Partie II --- Le delta de l'émetteur.}
\begin{enumerate}[resume]
\item Un mois avant l'échéance, à une volatilité de 25~\%, le spot étant supérieur de 1~\% à la barrière, donner le delta de l'émetteur par unité de notionnel.
\item Donner le delta après qu'une baisse de 1~\% a fait activer le knock-in de la note.
\item Quelle quantité du sous-jacent l'émetteur vend-il, pour 100 millions de notes~?
\item Donner le même résultat avec trois mois, puis avec un an restant.
\item Qu'a-t-il fait lorsque le spot est passé de 100~\% à 76~\%~?
\end{enumerate}

\medskip\noindent\textbf{Partie III --- Le marché.}
\begin{enumerate}[resume]
\item Si dix émetteurs ont émis des notes similaires avec la même barrière, que voit le marché~?
\item Quel ordre de grandeur Fang indique-t-il pour la couverture libérée dans la zone de la barrière~?
\item Pourquoi les contrats à terme sur l'indice sont-ils le lieu où le flux apparaît en premier~?
\item Quelle politique Lim et Choi ont-ils suggérée~?
\item Quelles données un régulateur aurait-il besoin de suivre pour surveiller ce risque~?
\end{enumerate}

\medskip\noindent\textbf{Partie IV --- Jugement.}
\begin{enumerate}[resume]
\item Comment intégrer dans le prix de la note le coût de liquidité de la falaise~?
\item Comment couvrir la position en amont de la barrière pour réduire la vente au moment du knock-in~?
\item Qui supporte la perte lorsque la barrière est franchie~?
\item Énoncer le \emph{résultat nommé}~: la quantité de sous-jacent que l'émetteur vend pour une baisse de 1~\% lorsque le spot est supérieur de 1~\% à la \omterm{def:dv:autocallables:protection}{barrière de
protection}, un mois avant l'échéance, pour 100 millions de notes.
\item En une phrase~: qu'est-ce qui fait d'une barrière de knock-in un événement de marché plutôt qu'un événement propre à une seule transaction~?
\end{enumerate}
\end{problem}

\section{Questions d'entretien}

\begin{interviewq}[$\star$ \iqroles{trader, banque}]\label{iq:dv:autocallables:1}
Décomposer un \omterm{def:dv:autocallables:autocallable}{autocallable} en options. Pourquoi peut-il verser un coupon élevé~?
\end{interviewq}

\begin{interviewq}[$\star\star$ \iqroles{chercheur}]\label{iq:dv:autocallables:2}
Avec quel modèle priceriez-vous un \omterm{def:dv:autocallables:autocallable}{autocallable}, et que néglige une volatilité plate~?
\end{interviewq}

\begin{interviewq}[$\star\star$ \iqroles{trader}]\label{iq:dv:autocallables:3}
Quelle est l'exposition de l'émetteur d'un \omterm{def:dv:autocallables:autocallable}{autocallable} au \omterm{def:dv:greeks-and-the-hedging-pnl:greeks}{vega}, au skew et au dividende, et comment celle-ci est-elle recyclée~?
\end{interviewq}

\begin{interviewq}[$\star\star$ \iqroles{développeur}]\label{iq:dv:autocallables:4}
Vos \omterm{def:dv:greeks-and-the-hedging-pnl:greeks}{grecques} Monte-Carlo pour un \omterm{def:dv:autocallables:autocallable}{autocallable} sont bruitées. Que faites-vous~?
\end{interviewq}

\begin{interviewq}[$\star\star$ \iqroles{risque}]\label{iq:dv:autocallables:5}
Comment mesureriez-vous le flux de couverture à l'échelle du marché que peut générer un ensemble de barrières de knock-in~?
\end{interviewq}

\begin{interviewq}[$\star\star\star$ \iqroles{trader, risque}]\label{iq:dv:autocallables:6}
L'indice est 2~\% au-dessus d'un niveau où une large part des notes du marché entrent en knock-in, un mois avant l'échéance de beaucoup
d'entre elles. Que faites-vous~?
\end{interviewq}
